\begin{figure*}[htbp!]
  \centering
  \includegraphics[width=\textwidth]{figure/compile.pdf}
  \caption{An overall example for compiling Hector to hardware. The original module with hybrid manners is split into two modules with dynamic and pipeline manners separately, and the connection between them is represented as a function call. Through resource and register sharing, pure \tor modules are transformed to \hec components containing compute units and registers. RTL implementations are generated by mapping ports and constructing controllers with the same functionality.}
  \label{fig:compile}
  \vspace{-9pt}
\end{figure*}

\section{Compiling IRs to Hardware}
\label{sec:compile}
%With the help of MLIR infrastructure, basic optimizations like dead code elimination (DCE) and common sub-expression elimination (CSE) can be easily implemented. Apart from these basic optimizations, hardware generation in Hector should go through these transformation passes:

To compile IR to hardware, Hector involves the following passes, Hybrid Time Graph Transformation (\autoref{sec:split}), Lowering from \tor to \hec (\autoref{sec:lower}), and RTL Generation (\autoref{sec:rtl}). \autoref{fig:compile} illustrates the main steps from \hec to RTL implementation written in Chisel. 
% \hl{can we change figure 4 to explain all the three steps ?}

The first pass converts the hybrid time graph in \tor into multiple pure time graphs by outlining connected sub-graphs. The second pass produces \hec programs with the same functionality of \tor through pipeline partition, control logic generation, and resource sharing. The third pass generates Chisel programs by implementing hardware controller and inserting necessary signals.

% \hl{we need a figure to explain this pass: inputs, outputs, transformations, etc. }

% \subsection{Verify Functionality}
% \label{sec:verify}
% The high-level topology enables early verification \hl{verification of what ? this needs examples} at a high level of abstraction. Some basic requirements of hardware design can be ensured by traversing the time graph. 

% \begin{itemize}
%     \item \textbf{Latency}: Because one variable cannot be used until the computation is complete, each path between its definition and usage should be shorter than the operation latency.
%     \item \textbf{Timing}: The total length of each combinational path cannot exceed the clock period, which can be examined by scanning each time edge.
%     \item \textbf{Memory Port}: The number of accesses in a single edge cannot be greater than the product of port limitation and edge length due to memory port limitations.
% \end{itemize}

\subsection{Hybrid Time Graph Transformation}
\label{sec:split}
\tor's semantics allows for the existence of multiple manners in the same time edge. However, because different types of hardware cannot communicate directly with one another, this hybrid time graph is unsuitable for subsequent transformations. This pass converts the hybrid time graph into multiple pure time graphs where all edges and nodes have the same manner. The connected sub-graph and related computation are defined as an external module, and communication between the outer and inner modules is accomplished through a function call. 

We calculate all live-in variables and live-out variables of connected sub-graph. Live-in variables are the arguments of new function, and that function return the values of all live-out variables. As shown in \autoref{fig:compile}(b), the connected sub-graph of hybrid time graph is marked by red rectangle. That hybrid time graph is split into two sub-graphs: a dynamic module and a pipeline module. The inner pipeline module take \verb|%c| and \verb|%d| as inputs, and return the value of \verb|%f|. This pass ensures that only one type of hardware is considered in the following transformations. 
% \hl{this subsection needs more details.}


\subsection{Lowering \tor to \hec}
\label{sec:lower}
The time graph in static modules is converted to a low-level FSM with the same functionality. For pipeline \tor modules, this pass traverses the time graph and calculates the "distance" of each node from the source node. The "distance" is indicated by “\#cycle” attributes of every time edge composing the path from the source node to the current node. The time nodes with the same "distance" will be placed in one stage. Each stage corresponds to a clock cycle, on which the allocated time node starts the execution of the operations whose start time is set as the current node. 

In contrast to \tor, \hec specifies where the computation takes place and where the intermediate value is stored. This opens up the possibility of resource and register sharing, which allows disjoint computations and values to share the same unit. \tor fits well in such optimizations because of the high-level control flow and timing information of computation contained in the time graph.

\textbf{Resource and register sharing}. The resource sharing can only happen between two operations that never occur at the same time, so two disjoint operations mustn't overlap in the time graph. We build a conflict graph that shows all potential conflicts and uses a coloring strategy to assign compute units to each node. As for register sharing, we use a live range analysis on the time graph to calculate the live range of each variables. The value of a variable in its live range must be stored in a register unless it is used in time. The main difference between register sharing and resource sharing is that each edge in the conflict graph represents an overlapping live range. 
% \hl{in 4.1, we mention a liveness analysis, is it the same ?}

As for hardware with dynamic behavior, each pair of operations may happen at the same time, restricting resource sharing. As a result, the lower pass just considers the generation of control logic. The handshake protocol, which describes data transfer with valid and ready signals, is used to implement the dynamic approach. This protocol is hidden in the \hec IR, which further simplifies the description. 
% \hl{we do not introduce the elastic components before.}
Elastic components~\cite{elastic_circuit} are adopted to solve the synchronization of tokens in the hardware with dynamic behavior. For example, \verb|branch| transfers the data to one of the outputs based on a condition, and \verb|merge| accepts one of the inputs which is similar to a $\phi$-function in software IR. We calculate the live-in and live-out variables of each regions, and a template-based approach is adopted to generate data transfer of variables among different regions. After that, \verb|fork| is inserted to synchronize multiple uses of the same value, which is shown in \autoref{fig:compile}(c). Both \verb|addf| and \verb|subf| need the value of \verb|%a|, so a \verb|fork| component is inserted for that variable. The two outputs of \verb|fork| are fed into \verb|addf| and \verb|subf| separately.

\subsection{RTL Generation}
\label{sec:rtl}
\hec contains sufficient information to generate synthesizable RTL. Each component in \hec has explicit I/O port definitions, assignments of ports, a controller, and allocations of register and resource (including memory, compute units, and other sub-component). The RTL generation pass creates a Chisel \cite{chisel} program by mapping each component to a Chisel module with the same ports and assigning the corresponding value to these ports. Chisel will automatically insert the multiplexer, allowing for multiple assignments under different conditions. All of the built-in resources, such as memory and compute units are implemented as a pre-defined Chisel Module. 

\textbf{Controller Generation}. The STG-style component is controlled by a state transition graph that contains a state set and the transition between them, whereas the pipeline-style component is controlled by a collection of stages and matching II. By inserting some extra registers and wires, this pass builds a Chisel-style FSM for these two types of components with the same functionality in \tor. Due to the need to resolve conflicts at run-time, handshake-style components do not have a centralized controller. For each port, the handshake protocol is implemented as a Chisel class \texttt{DecoupledIO} with ready, valid, and data signals. The Chisel program for a pipeline component and a dynamic component is shown in \autoref{fig:compile}. The controller of the pipeline component is eliminated because the II equals one.

\textbf{Wrapper construction}. Due to the distinction between handshake-style component and the other two types of components, this pass additionally generates a wrapper that attaches the handshake protocol to naked components. Whether an STG-style (or pipeline-style) component is valid and ready is determined by the component's completion (or the validity of the initial stage and final stage).